% -------------------------------------------------------------- %
%      csd5748 - HY-225 - Spring 2025 - CSD-UOC - Ex. 14      %
% -------------------------------------------------------------- %

\documentclass{article}

\usepackage[utf8]{inputenc} 
\usepackage[english, greek]{babel}
\usepackage{subcaption}
\usepackage{hyperref}
\usepackage{xcolor}
\hypersetup{ hidelinks, colorlinks=true, linkcolor=blue, urlcolor=blue }
\captionsetup[subfigure]{labelformat=empty}
\usepackage{amsmath}
\usepackage{amsfonts}
\usepackage{amssymb}
\usepackage{array}
\usepackage{graphicx}
\usepackage[a4paper, left=0.8cm, right=1.4cm, top=2cm, bottom=2cm]{geometry}
\usepackage[absolute, overlay]{textpos}
\usepackage{placeins}
\usepackage{listings}
\usepackage{xcolor}
\setlength{\jot}{0pt}
\setlength{\mathsurround}{0.1cm}
\setlength{\jot}{0.1cm}
\newcommand{\skato}{\hspace{0.25cm}}
\newcommand{\skata}{\hspace{1cm}}
\usepackage[english, greek]{babel}  
\usepackage{tcolorbox} 
\usepackage{titlesec}
\usepackage{fancyvrb} 
\geometry{margin=2.5cm}
\lstset{
  language=C,
  inputencoding=utf8,
  extendedchars=true,
  basicstyle=\ttfamily\small,
  keywordstyle=\color{blue},
  commentstyle=\color{gray}\itshape,
  stringstyle=\color{red},
  showstringspaces=false,
  columns=fullflexible,
  keepspaces=true,
  frame=single,
  breaklines=true,
  tabsize=2,
  escapeinside={(*@}{@*)}, 
  morekeywords={char, return, if, while, volatile, unsigned, int, define},
}

\begin{document}

% -------------------------------------------------------------- %

\text{   }
\vspace{-1.0cm}
\begin{center}
\setlength{\fboxrule}{2pt}
\setlength{\fboxsep}{15pt} 
\fbox
{
{\noindent{\Huge \textbf{\textlatin{HY225 -} \textbf{Ασκήσεις 14} \textlatin{- csd5748}}}}
}
\end{center}
\vspace{2.000cm}


% -------------------------------------------------------------- %


\begin{center}
    {\underline{\Large \textbf{Άσκηση 14.1}}}
\end{center}
\vspace{0.1cm}
 

\setlength{\fboxrule}{1pt}
\begin{figure}[h]
  \centering
  \fbox{\includegraphics[width=0.55\textwidth]{14ai.png}}
  \caption*{ }
\end{figure}
 


\vspace{0.1cm} 

\noindent \boxed{\text{α}} Όταν η συσκευή Ε/Ε στέλνει δεδομένα στη μνήμη μέσω \textlatin{DMA}, ο επεξεργαστής διαβάζει τα σωστά αν γίνει \textlatin{cache miss} (οπότε προσπελάζει την κύρια μνήμη) ή αν υπάρχει μηχανισμός που ενημερώνει την κρυφή μνήμη. Διαφορετικά, σε \textlatin{cache hit} χωρίς ενημέρωση, διαβάζει παλιά και λανθασμένα δεδομένα ...

\vspace{0.25cm}

\noindent \boxed{\text{β}} Όχι, γιατί με πολιτική \textlatin{write-through} κάθε εντολή \textlatin{store} ενημερώνει αμέσως τη μνήμη. Άρα, όταν ξεκινά το \textlatin{DMA}, η μνήμη έχει ήδη τα σωστά δεδομένα για τη συσκευή εξόδου ...

\vspace{0.25cm}

\noindent \boxed{\text{γ}} Με \textlatin{write-back}, αν η έξοδος διαβάσει από τη μνήμη πριν τα δεδομένα περάσουν από την \textlatin{cache}, ενδέχεται να είναι λανθασμένα. Αν έχουν ήδη γραφτεί τότε θα είναι σωστά ...



% ---------------------------------------------------------------------------------------



\newpage

\begin{center}
    {\underline{\Large \textbf{Άσκηση 14.3}}}
\end{center}

\bigskip
\bigskip
\bigskip


\noindent \boxed{\text{α}} 
\textlatin{ii}: 1 \textlatin{read}: \text{Η κρυφή μνήμη του} \textlatin{A} \text{έχει αστοχία οπότε ζητά την μεταβλητή από την κύρια μνήμη.}

\noindent \boxed{\text{α}} 
\textlatin{iii}: 1 \textlatin{read}:
\text{Η κρυφή μνήμη του} \textlatin{B} \text{έχει αστοχία οπότε ζητά την μεταβλητή από την κύρια μνήμη.}

\noindent \boxed{\text{α}} 
\textlatin{iv}: \text{Καμιά πράξη}: \text{Η κρυφή μνήμη} \textlatin{A} \text{ενημερώνεται με την τιμή 11 αλλά όχι η κύρια μνήμη \textlatin{(write-back)}}.

\noindent \boxed{\text{α}} 
\textlatin{v}: \text{Καμιά πράξη}: \text{Η κρυφή μνήμη} \textlatin{B} \text{έχει ευστοχία (χρησιμοποιεί παλιά τιμή 0), οπότε δεν ζητά τίποτα.}


\noindent \boxed{\text{α}} 
$*$: \textlatin{write-through}: (\textlatin{iii}) \text{Κρυφή μνήμη Α ενημερώνει τιμή 11 παντού (\textlatin{iv}) άρα την χρησιμοποιεί ο Β ...}


\bigskip
\bigskip

\noindent \boxed{\text{β}} 
\textlatin{ii}: 1 \textlatin{read}: \text{Κρυφή μνήμη Α αστοχία, ζητάει από κύρια μνήμη μέσω αρτηρίας και κατάσταση γίνεται \textlatin{S}.}

\noindent \boxed{\text{β}} 
\textlatin{iii}: 1 \textlatin{read}: \text{Κρυφή μνήμη Β αστοχία, ζητά μεταβλητή από κρυφή μνήμη Α και κατάσταση γίνεται \textlatin{S}.}

\noindent \boxed{\text{β}} 
\textlatin{iv}: 1 \textlatin{inv} + 1 \textlatin{wr}: \text{Α γράφει 11 γίνεται (Μ),} \textlatin{invalidate} \text{\textlatin{bus}, κρυφή Β ακυρώνεται (\textlatin{I}), κύρια ενημερώνεται.}

\noindent \boxed{\text{β}} 
\textlatin{v}: 1 \textlatin{read}: \text{Κρυφή του} \textlatin{B} \text{έχει αστοχία (}\textlatin{I}\text{), οπότε ζητά από τον Α} \text{(}\textlatin{M}\text{) και παίρνει 11*, κατάσταση \textlatin{B (S)}.}

\noindent \boxed{\text{β}} 
$*$: \ 
\text{Με κατάσταση Ε}: \ 
\text{(\textlatin{ii}) Α (Ε) αντί (\textlatin{S}) αφού μόνος με αντίγραφο. (\textlatin{iii}) Α (Ε) σε (\textlatin{S}) λόγω Β ...}



\bigskip
\bigskip

\noindent \boxed{\text{γ}} 
\textlatin{ii}: 1 \textlatin{read}: \text{Κρυφή μνήμη} \textlatin{A} \text{αστοχία, ζητά μετ. από κύρια μνήμη, κατάσταση} \textlatin{S} \text{(πρωτόκολο συνοχής \textlatin{MSI}).}

\noindent \boxed{\text{γ}} 
\textlatin{iv}: 1 \textlatin{inv} + 1 \textlatin{wr}: \textlatin{A} \text{γράφει 11, κατάσταση} \textlatin{M}\text{,} \textlatin{invalidate} \text{στο} \textlatin{bus} \text{(χωρίς λόγο), κύρια μνήμη ενημερώνεται.} 


\noindent \boxed{\text{γ}} 
\textlatin{ii}: 1 \textlatin{read}: \text{Κρυφή μνήμη} \textlatin{A} \text{αστοχία, ζητά μεταβλητή από κύρια μνήμη, κατάσταση} \textlatin{E} \text{(μοναδικός)(\textlatin{MESI}).}

\noindent \boxed{\text{γ}} 
\textlatin{iv}: 1 \textlatin{wr}: \textlatin{A} \text{γράφει 11, κατάσταση} \textlatin{M}\text{,} \text{χωρίς} \textlatin{invalidate} \text{(} \textlatin{B} \text{δεν έχει αντίγραφο), κύρια μνήμη ενημερώνεται.} 






% ---------------------------------------------------------------------------------------



\newpage

\begin{center}
    {\underline{\Large \textbf{Άσκηση 14.4}}}
\end{center}

\bigskip
\bigskip
\bigskip



\noindent \boxed{\text{α}} 
{\textlatin{Assembly risc-v with no instruction scheduling ...}}

\selectlanguage{english}

\begin{lstlisting} 
                      * lw t0, 120          # Load a from DM[120] to t0.
                      * addi t1, t0, 2      # t1 = t0 + 2 (x = a + 2).
                      * sw t1, 128          # Store t1 (x) at DM[128].
                      * lw t2, 124          # Load b from DM[124] to t2.
                      * addi t3, t2, -1     # t3 = t2 - 1 (y = b - 1).
                      * sw t3, 132          # Store t3 (y) at DM[132]. 
\end{lstlisting}

\selectlanguage{greek}

\vspace{-0.1cm}
\noindent \boxed{*} \
\textlatin{Pipeline} \textlatin{(IF, ID, EX, MEM, WB)}: \ \textlatin{Load} \& \textlatin{Addi}
$\rightarrow$ 2 \text{κύκλοι} $\Rightarrow$ \text{Χαμένοι κύκλοι} = 2$\times$2 = 4 ...


\bigskip
\bigskip
\bigskip


\noindent \boxed{\text{β}} 
{\textlatin{Assembly risc-v with instruction scheduling ...}}
\vspace{0.1cm}

\selectlanguage{english}

\begin{lstlisting} 
                      * lw t0, 120          # Load a from DM[120] to t0.
                      * lw t2, 124          # Load b from DM[124] to t2.
                      * addi t1, t0, 2      # t1 = t0 + 2 (x = a + 2).
                      * addi t3, t2, -1     # t3 = t2 - 1 (y = b - 1).
                      * sw t1, 128          # Store t1 (x) at DM[128].
                      * sw t3, 132          # Store t3 (y) at DM[132].
\end{lstlisting}

\selectlanguage{greek}


\bigskip
\bigskip
\bigskip


\noindent \boxed{\text{γ}} 
{\textlatin{Assembly risc-v with no pointers instruction scheduling ...}}
\vspace{0.1cm}

\selectlanguage{english}

\begin{lstlisting} 
                      * lw t0, 0(x14)       # Load *pa from x14 to t0.
                      * addi t1, t0, 2      # t1 = t0 + 2 ((*px)=(*pa)+2).
                      * sw t1, 0(x16)       # Store t1 to x16 (*px).
                      * lw t2, 0(x15)       # Load *pb from x15 to t2.
                      * addi t3, t2, -1     # t3 = t2 - 1 ((*py)=(*pb)-1).
                      * sw t3, 0(x17)       # Store t3 to x17 (*py).
\end{lstlisting}

\selectlanguage{greek}

\vspace{-0.1cm}

\noindent \boxed{*} \
\text{Ο \textlatin{compiler} δεν μπορεί να κάνει \textlatin{instruction scheduling} αφού δεν γνωρίζει τις τιμές των \textlatin{pointers} ...}

\noindent \boxed{*} \
\text{Εάν} \textlatin{*px}$\neq$\textlatin{*pb} \text{τότε δεν υπάρχουν αλληλεξαρτήσεις} $\Rightarrow$ \text{μπορεί να υλοποιηθεί \textlatin{instruction scheduling} ...}

\noindent \boxed{*} \
\text{Εάν} \textlatin{*px}$=$\textlatin{*pb} \text{τότε υπάρχουν αλληλεξαρτήσεις} $\Rightarrow$ \text{απαγορεύεται υλοποίηση \textlatin{instruction scheduling} ...}


\bigskip 
\bigskip 
\bigskip


\noindent \boxed{\text{δ}}
\text{Η δεύτερη εντολή μπορεί να εξαρτάται από την πρώτη αν οι δείκτες δείχνουν στην ίδια διεύθυνση (1\%).}

\vspace{0.1em}

\noindent \boxed{*}
\text{Στην αντίθετη περίπτωση που οι διευθύνσεις είναι διαφορετικές δεν υπάρχουν αλληλεξαρτήσεις (99\%).} 

\noindent \boxed{*}
\text{Η επόμενη \textlatin{load} μπορεί να εξαρτάται από την προηγούμενη \textlatin{store}, υπό την προϋπόθεση} \textlatin{*px}$=$\textlatin{*pb} ...

\noindent \boxed{*}
\text{Έτσι, \textlatin{hardware} πετυχαίνει δυναμική αναδιάταξη και καθυστερεί δεύτερη μέχρι ολοκληρωθεί πρώτη ...}


\noindent \boxed{*}
\text{Τελικά, χρόνος χάνεται όταν βρισκόμαστε στο 1\% των περιπτώσεων και όχι στο 99\% ...}






% ---------------------------------------------------------------------------------------



\newpage

\begin{center}
    {\underline{\Large \textbf{Άσκηση 14.5}}}
\end{center}

\bigskip
\bigskip
\bigskip


\noindent \boxed{\text{α}} \
\text{(Συνολικοί κύκλοι αναμονής)}$-$\text{(Χρήσιμες εντολές που εκτελούνται)} $=$ $80 - 8$ $=$ 72 \text{κύκλοι.}
 
 
\bigskip

\noindent \boxed{\text{β}} \
\text{Το νήμα (πρόγραμμα) A τρέχει σε ακριβώς τον ίδιο χρόνο, αλλά η συνολική απόδωση βελτιώνεται ...}

 
\bigskip

\noindent \boxed{\text{γ}} \
\text{Οι συνολικά χαμένοι κύκλοι είναι 0, ο επεξεργαστής δεν αδρανοποιείται, πάντα εκτελείται κάτι χρήσιμο.}



 


% ---------------------------------------------------------------------------------------

\newpage

\begin{textblock}{10}(14.401, 15.501)
    \boxed{\textbf{\LaTeX{}}}
\end{textblock}

\begin{center}
    \vspace*{\fill}
    \textbf{\LARGE{ΤΕΛΟΣ}} \\[0.001cm]
    \textbf{\hspace{0.1cm}(ευχαριστώ)}
    \\[0.1cm]
    \vspace*{\fill}
\end{center}

% ---------------------------------------------------------------------------------------

\end{document}

% -------------------------------------------------------------- %
%                           The End :)                           %
% -------------------------------------------------------------- %

